\documentclass[10pt,a4paper]{amsart}
\usepackage[margin=19mm]{geometry}
\usepackage{amsmath,amssymb,mathtools}
\usepackage[scr=rsfs,cal=euler]{mathalfa}
\usepackage{xfrac}
\usepackage[unicode,hidelinks,bookmarks=false]{hyperref}
\newcommand{\ie}{i.e.\ }
\newcommand{\cf}{cf.\ }
\newcommand{\resp}{respectively\ }
\newcommand{\Subsection}{\subsection}
\providecommand{\nobreakdash}{\nobreak\hbox{-}}
\setlength{\emergencystretch}{2em}
\multlinegap0pt
\usepackage{amssymb}
\usepackage{enumerate}
\usepackage{mathtools}
\newtheorem{lem}{Lemma}
\newtheorem{definition}{Definition}
\newtheorem{thm}{Theorem}
\newcommand{\R}{\mathbb{R}}
\DeclareMathOperator*{\arginf}{arg\,inf}
\newcommand{\loc}{\mathrm{loc}}
\title{Regularity of a free boundary determined by the value of the gradient.\\ Part 1: Third order asymptotics.}
\author{Frida Fejne}
\date{Source: arXiv:2609.11324v1 (CC BY 4.0)}
\begin{document}
\maketitle

\noindent\emph{Source and licence.} Regularity of a free boundary determined by the value of the gradient.\\ Part 1: Third order asymptotics.. Version arXiv:2609.11324v1 (CC BY 4.0), distributed by arXiv under the Creative Commons Attribution 4.0 International licence, http://creativecommons.org/licenses/by/4.0/, as recorded in the arXiv licence metadata for this exact version and shown on the abstract page https://arxiv.org/abs/2609.11324v1. Full licence notice: \texttt{licenses/arxiv-2609.11324-CC-BY-4.0.txt}.\par\smallskip

\noindent\textit{Editorial scope.} Counted unit: the article's thm thm:mainthm (source0.tex lines 2704--2716). The article's complete original proof of it is delivered with it (source0.tex lines 2717--2859, one \texttt{proof} environment of 143 lines). No mathematics is written, repaired or re-derived: the statement and the proof are the article's own lines, in the article's own notation, and no retained line is substituted or re-flowed.  Retained uncounted as the source-local context the proof needs: lines 77--80; lines 608--634; lines 1335--1379; lines 1949--1964; lines 2187--2210; lines 2193--2195; lines 2201--2203; lines 2369--2409; lines 2393--2396; lines 2462--2464; lines 2470--2478; lines 2679--2681; lines 2683--2685.  Nothing was omitted from any retained span, so the package carries no omission record.  No retained line contains a citation, so the wrapper carries no bibliography and no bibliography entry.  No retained line contains a figure environment, an image-inclusion command or drawing code, so no image is delivered and the package declares no asset dependency.  Editorial formatting only, in addition: an AMS \texttt{amsart} preamble that restates the article's own declarations and macros used by the retained lines, namely the environments \texttt{lem}, \texttt{definition}, \texttt{thm} on separate counters, together with the standard packages those lines need.  The retained environments print in this document as Lemma 1, Definition 1, Theorem 1 in order of appearance, whereas the article prints the same items with the numbering of its own full document.  The complete native source of the article as distributed in the arXiv e-print for this exact version is bundled unchanged as \texttt{sources/210008/source0.tex}.
\begin{equation}
\label{eq:varint}
    J(u) : = \int_{\Omega} F(\nabla u) \ dx,
\end{equation}

\begin{lem}\label{cor:C3result}
 Let $m(x')$ be a second degree homogeneous polynomial, 
 $$
 m(x')= \sum_{i,k=1}^{n-1}m_{ik} x_i x_k,
 $$
 $M \in \R$ a constant and
 $$
 r(x)=\left\{ \begin{array}{ll}
               l_1\cdot x & \textrm{for }x_n\ge 0 \\
               l_2\cdot x & \textrm{for }x_n<0
              \end{array}
 \right.
 $$
 be a function that is linear in $B_1^+$ and $B_1^-$.
Assume that $q\in W^{1,2}_{\textup{loc}}(B_{1})$ is a solution to the following PDE:
 \begin{align}\label{eq:lin_pde_withmr}
 0 &= \int_{B_{1}^+}\sum_{i,j=1}^{n}F^1_{p_ip_j}(\nu) D_i qD_j\varphi dx+
 \int_{B_{1}^-}\sum_{i,j=1}^{n}F^2_{p_ip_j}(\nu) D_i qD_j\varphi dx \\
 & \qquad +\int_{B_1}r(x)\varphi(x) dx +\int_{B_1'}(M+m(x'))\varphi(x') dx', \nonumber
 \end{align}
 where $\varphi \in C^{\infty}_0(B_{1})$ and $F^m_{p_i p_j}(\nu)$ are constants for every $m=1,2$ and $i,j=1, \hdots n$. Then $q$ will be $C^3_{\loc}(B_1^{\pm})$, and for any $r \in (0,1)$ there is a
 constant $C$ such that
 $$
 \|D^3 q\|_{L^{\infty}(B^+_{r})}+\|D^3 q\|_{L^{\infty}(B^-_{r})}\leq C \left( \|q\|_{L^2(B_1)}+ \|r\|_{L^2(B_1)}+ \|m\|_{L^2(B_1')}\right),
 $$
 where $C$ only depends on $r$, $F^1$ and $F^2$.
\end{lem}

\begin{lem}\label{lem:structure_p}
Let $p \in \Pi(\nu)$ where $\nu = (\nu_1,\nu_2, \hdots, \nu_n)$ is a unit vector. Furthermore, let $J$ denote the jump in the second derivatives of $F$ at $\nu$, and in direction $\nu$, that is,
\begin{equation*}\label{eq:DefJump}
J=\nu\cdot\big(D^2 F^1(\nu)-D^2 F^2(\nu)\big) \nu^T.
\end{equation*}
Then the following holds:
\begin{enumerate}
 \item when not both $i=n$ and $j=n$ we have that
 $$
D^2_{ij} p_- = D^2_{ij} p_+,
 $$
 \item we have that
 \begin{equation*}
     \sum_{k=1}^{n} \nu_k D^2_{mk} p = 0,
 \end{equation*}
  for $m=1, \hdots, n-1$, and 
 \begin{equation*}
    \nu \cdot \nabla p(x) = D^2_{\nu n} p  x_n = \sum_{k=1}^n \nu_k D^2_{nk}p x_n,
\end{equation*}
 \item the second derivatives of $p_\pm$ are related by
\begin{equation}\label{eq:Dpn1}
  D^2_{nn} p_-=D^2_{nn} p_+
 +\frac{J \nu_n D^2_{\nu n}p_-}{F^1_{p_n p_n}(\nu)},   
\end{equation}
 \item and the first derivatives of $p_{\pm}$ are related by
 $$
 \nu\cdot \nabla p_-=\nu\cdot \nabla p_+ +\frac{JD^2_{\nu n}p_-}{F^1_{p_n p_n}(\nu)} (\nu_n)^2 x_n,
 $$
 \item the following identities are true:
    \begin{align}
        \sum_{i, k = 1}^n F^1_{p_i p_k}(\nu)D^2_{ik} p_- &= (\nu \cdot (D^2 F^1 (\nu)-D^2 F^2(\nu)) \nu^T) \nu_n D^2_{\nu n}p_-, \label{eq:sumpoly1} \\
         \sum_{i, k = 1}^n F^2_{p_i p_k}(\nu)D^2_{ik} p_+ & = -(\nu \cdot (D^2 F^1 (\nu)-D^2 F^2(\nu)) \nu^T) \nu_n D^2_{\nu n}p_+, \label{eq:sumpoly2}
    \end{align}
\item $D^2_{\nu n}p_{+}=0$ if and only if $D^2_{\nu n}p_{-}=0$.
 \item if $D^2_{\nu n}p_{-} \neq D^2_{\nu n}p_{+}$, then $\nu_n \neq 0$ and $J \neq 0$. Furthermore,
        \begin{equation*}
                D^2_{\nu n} p_- F^2 _{p_n p_n}(\nu) = D^2_{\nu n } p_+ F^1_{p_n p_n}(\nu),
        \end{equation*}
and
$$
F^1_{p_n p_n}(\nu) \neq F^2_{p_n p_n}(\nu).
$$
\end{enumerate}

\end{lem}

\begin{definition}
\label{def:pk}
{\sl For a minimizer $u$ to \eqref{eq:varint} in $B_1$ we define
\begin{equation*}
    a, \nu, \hat{p} = \arginf_{a + \nu \cdot x + p \in \mathcal{M}(\nu)} \left\|u(x)-a -\nu \cdot x - p(x) \right\|_{L^2(B_1)}.
\end{equation*}
 Furthermore, we let $\delta=  \|\hat{p} \|_{L^2(B_1)}$ and $p = \hat{p}/\delta$, and define
\begin{align*}
    \epsilon  &=  \frac{1}{\delta}  \left\|u(x) -a - \nu \cdot x - \delta p(x) \right\|_{L^2(B_1)}.
\end{align*}
Assuming that $\epsilon \neq 0$ we define
\begin{align*}
     q &= \frac{1}{\epsilon \delta} \left(u(x) -a - \nu \cdot x - \delta p(x) \right),
\end{align*}
so that $\left\|q(x) \right\|_{L^2(B_1)} = 1$.}
\end{definition}

\begin{lem}\label{lem:summaryuj}
   From assumptions 1-4 in the beginning of this section it follows that each $u^j$ can be expressed as
$$
    u^j=a^j+\nu^j\cdot x+\delta_jp^j+\epsilon_j\delta_j q^j,
$$
where $a^j,\nu^j,p^j,q^j,\epsilon_j$ and $\delta_j$ are as in Definition \ref{def:pk}. Furthermore, $\epsilon_j \to 0$, $\delta_j \to 0$, $\nu^j \to \nu$, and $p^j \to p^0$ in $C^{1, \alpha}(B_1)$ as $j \to \infty$. In addition, we have that
    \begin{equation}\label{eq:qjass1}
        \epsilon_j q^j \to 0 \quad \textup{in} \quad C^{1,\alpha}_{\loc}(B_1),
    \end{equation}
and
\begin{equation}\label{eq:qjass2}
  \frac{|\nu^j|-1}{\delta_j} \leq \|\epsilon_jq^j\|_{C^{1,\alpha}(B_{1/2})}\to 0,
 \end{equation}
as $j \to \infty$. Moreover,
    \begin{equation}\label{eq:qjass3}
 \nu_n \neq 0 \quad \textup{and} \quad \left|\nu \cdot [D^2 F^1(\nu)- D^2 F^2(\nu)]  \nu^T \right| \geq C,
\end{equation}
and there exists some constant $c_{ndg} >0$ such that 
 \begin{equation}\label{eq:qjass4}
     c_{\nu^j}^{\pm} \geq c_{ndg},
 \end{equation}
for sufficiently large values of $j$.

\end{lem}

    \begin{equation}
        \epsilon_j q^j \to 0 \quad \textup{in} \quad C^{1,\alpha}_{\loc}(B_1),
    \end{equation}

    \begin{equation}
 \nu_n \neq 0 \quad \textup{and} \quad \left|\nu \cdot [D^2 F^1(\nu)- D^2 F^2(\nu)]  \nu^T \right| \geq C,
\end{equation}

\begin{lem}\label{lem:conv_AjikCmjikl}
Let $0<r<1$ and $u^j$ be as in Lemma \ref{lem:summaryuj}. Then the following holds:
 \begin{enumerate}
     \item the coefficients $E^{m,j}_{ikl}$ can be estimated as
 \begin{equation}\label{eq:Eikl}
   |E^{m,j}_{ikl}| \leq C \delta_j  \|\epsilon_j q^j\|_{C^{1, \alpha}(B_{1/2})},
 \end{equation}
 for $m=1,2$,
 \item for $x \in B_r$ we also have that
\begin{equation}
\label{eq:Ccoeffconv}
       \lim_{j\to \infty} C^{m,j}_{ikl}(x) =  F^m_{p_i p_k p_l}(\nu)\textup{ for } m = 1,2,
\end{equation}
where the convergence is uniform in $B_r^{\pm}$,
\item 
there exists a function $g^r_j$ that is uniformly bounded for each $j$:
$$
g_r^j = \max \!\left(\left|a^{j,+}_r\right|, \left|a^{j, -}_r+ \frac{\delta_j \|\nabla (p^j +\epsilon_j q^j)\|^2_{C^{\alpha}(B_r)}}{2c_{\nu^j}^{-}}\right| \right) \!,
$$
where
$$
a^{j, \pm}_r = \frac{|1-|\nu^j|^2|}{2\delta_j c_{\nu^j}^{\pm}}+\frac{ \epsilon_j \|\nabla q^j\|_{C^{\alpha}(B_r)}}{c_{\nu^j}^+},
$$
such that $g^j_r \to 0$ as $j\to \infty$, and such that the following holds:
\begin{align}\label{eq:fbestimate1}
    B_r& \cap \{x_n>g^j_r \} \subset \{ x \in B_1: |\nabla u^j(x)|>1 \}, \\
    B_r& \cap \{x_n < -g^j_r \} \subset \{ x \in B_1: |\nabla u^j(x)|<1 \}, \label{eq:fbestimate2}
\end{align}
\item
for $x \in B_r$ it follows that
\begin{equation}
\label{eq:Acoeffconv}
\lim_{j\to \infty} A^j_{ik}(x) =\left\{\begin{array}{ll}
                                        F^1_{p_ip_k}(\nu) & \textrm{if }x_n > 0, \\
                                        F^2_{p_ip_k}(\nu) & \textrm{if }x_n< 0,
                                       \end{array}
\right.
\end{equation}
where the convergence is uniform in $B_r^{\pm}$.
 \end{enumerate}
\end{lem}

\begin{align}
    B_r& \cap \{x_n>g^j_r \} \subset \{ x \in B_1: |\nabla u^j(x)|>1 \}, \\
    B_r& \cap \{x_n < -g^j_r \} \subset \{ x \in B_1: |\nabla u^j(x)|<1 \}, 
\end{align}

\begin{equation}\label{eq:ass1}
    \tau_j:= \frac{\delta_j}{\epsilon_j}\leq C_1,
\end{equation}

\begin{lem}
\label{lem:weak_convergence}
Let $u^j$ be as in Lemma \ref{lem:summaryuj} and assume that \eqref{eq:ass1} holds. Then, for large values of $j$, the sequence $q^j$ satisfies
\begin{align} \label{eq:qjBound}
 \int_{B_1} |\nabla q^j|^2 \xi^2 \ dx  & \leq C(|K_j|+1) \left(1 + \left(\int_{B_1} \xi^2 |\nabla \xi|^2 |q^j|^2 dx \right)^{1/2}+\beta \int_{B_1} \xi^2 |\nabla q^j|^2 dx + \frac{1}{4 \beta} \int_{B_1} \xi^2 dx \right)  \\
 &\quad + C \left(\int_{B_1} |\nabla \xi|^2 |q^j|^2 \ dx + \left(\int_{B_1} |\xi q^j|^2 dx \right)^{1/2} \left( \int_{B_1} \xi^2 dx \right)^{1/2} \right), \nonumber
\end{align}
for any $\beta >0.$
\end{lem}

\begin{equation}\label{eq:ass2}
    \frac{||\nu^j|^2-1| }{\delta_j \epsilon_j} \leq C_2.
\end{equation}

\begin{equation}\label{eq:alphatau}
    \tau = \lim_{j \to \infty} \frac{\delta_j}{\epsilon_j} \ \textup{ and } \ \alpha=\lim_{j\to \infty}  K_j= \frac{|\nu^j|^2-1}{\delta_j \epsilon_j}.
\end{equation}

\begin{thm}
\label{thm:mainthm}
 Let $u^j$ be as in Lemma \ref{lem:summaryuj}. Under assumptions \eqref{eq:ass1} and \eqref{eq:ass2} there is a subsequence of $q^j$, that we will continue to denote by $q^j$ such that $q^j\to q^0$ weakly in $ W_{\loc}^{1,2}(B_1)$ and satisfies
\begin{align}
\label{eq:q0equation}
     0 &= - \int_{B_1}  \sum_{i, k = 1}^n  F^*_{p_ip_k}(\nu,x) D_k q^0 D_i \varphi \, dx +  \tau \int_{B_1} \sum_{i, k, l = 1}^n   F^*_{p_ip_k p_l}(\nu,x) D_{ik} p^0 D_l p^0 \varphi \, dx \nonumber \\
      & \quad +  \frac{\nu \cdot (D^2 F^1 (\nu)-D^2 F^2(\nu)) \nu^T}{2} \nu_n \int_{B_1'}  \varphi(x',0) \left(\alpha + \tau |\nabla p^0(x',0)|^2\right) \, dx',
\end{align}
for any $\varphi \in C^{\infty}_0(B_1)$ and where $\alpha$ and $\tau$ are defined in \eqref{eq:alphatau}. Furthermore, for each $0<r<1$ there exists a constant $C_r$ such that
\begin{equation}\label{eq:C3est}
\|q^{0}\|_{C^{3}(B_r^{+})} < C_r \textrm{ and }\|q^{0}\|_{C^{3}(B_r^{-})} < C_r.
\end{equation}
\end{thm}

\begin{proof}
We take $0<r<1$, and recall that $q^j$ satisfies the PDE
\begin{align}
    0  &= -\int_{B_1}  \sum_{i, k = 1}^n  A^j_{ik}(x) D_k q^j D_i \varphi \ dx +\frac{1}{\epsilon_j}  \int_{B_1^+ \setminus D^j_+}  \sum_{i, k,l = 1}^n E^{1,j}_{ikl} \nu^j_l  D^2_{ik}p^j \varphi dx \nonumber\\
    &\quad +\frac{1}{\epsilon_j}   \int_{B_1^- \setminus D^j_-}  \sum_{i, k,l = 1}^n E^{2,j}_{ikl} \nu^j_l  D^2_{ik}p^j \varphi dx +\frac{\delta_j}{\epsilon_j} \int_{(B^+_1 \setminus D^j_+) \cup D^j_-} \sum_{i, k, l = 1}^n  C^{1,j}_{ikl}(x) D^2_{ik} p^j D_l p^j \varphi \ dx \nonumber \\
    &\quad + \frac{\delta_j}{\epsilon_j} \int_{(B^-_1 \setminus D^j_-) \cup D^j_+ } \sum_{i, k, l = 1}^n  C^{2,j}_{ikl}(x) D^2_{ik} p^j D_l p^j \varphi \ dx + \frac{1}{\epsilon_j} \int_{D^j_-}  \sum_{i, k = 1}^n   F^1_{p_i p_k}(\nu^j) D^2_{ik} p^j \varphi \ dx \nonumber \\
    &\quad + \frac{1}{\epsilon_j} \int_{D^j_+}  \sum_{i, k = 1}^n   F^2_{p_i p_k}(\nu^j) D^2_{ik} p^j \varphi \ dx, \label{eq:pdeforqj} 
\end{align}
for any $\varphi \in C^{1,2}(B_r) $. From Lemma \ref{lem:conv_AjikCmjikl} we know that 
\begin{equation}
\label{eq:Acoeffconv2}
\lim_{j\to \infty} A^j_{ik}(x) =\left\{\begin{array}{ll}
                                        F^1_{p_ip_k}(\nu) & \textrm{if }x_n > 0, \\
                                        F^2_{p_ip_k}(\nu) & \textrm{if }x_n< 0,
                                       \end{array}
\right.
\end{equation}
and the convergence is locally uniform in $B_1^+$ and $B_1^-$, respectively. We take $\lambda>0$. From \eqref{eq:fbestimate1} and \eqref{eq:fbestimate2} it follows that we can find $S_{\lambda}= \{x \in B_r: |x_n|\leq \lambda \}$ such that $\Gamma^j \subset S_\lambda$ for $j$ large enough. Thus, we can write 
\begin{align*}
&\int_{B_1}  \sum_{i, k = 1}^n  A^j_{ik}(x) D_k q^j D_i \varphi \ dx \\
     &= \int_{B_1^+\setminus S_{\lambda}}  \sum_{i, k = 1}^n  A^j_{ik}(x) D_k q^j D_i \varphi \ dx 
     + \int_{B_1^-\setminus S_{\lambda}}  \sum_{i, k = 1}^n  A^j_{ik}(x) D_k q^j D_i \varphi \ dx 
     + \int_{S_{\lambda}}  \sum_{i, k = 1}^n  A^j_{ik}(x) D_k q^j D_i \varphi \ dx \\
     &= I_1(j, \lambda) + I_2(j, \lambda) +I_3(j, \lambda).
\end{align*}
We note that since $\nabla F$ is Lipschitz we have that $\|D^2 F \|_{L^{\infty}(B_1)} \leq C$, and therefore 
$$
\|A^j_{ik}(x)\|_{L^{\infty}(B_1)} \leq C.
$$
We use \eqref{eq:Acoeffconv2}, Lemma \ref{lem:weak_convergence}, and the fact that the product of a uniformly and weakly convergent sequence is weakly convergent, to conclude that
\begin{align*}
     I_1(j, \lambda) &\to I_1(\lambda)  = \int_{B_1^+ \setminus {S_{\lambda}}} \sum_{i, k = 1}^n  F^1_{p_ip_k}(\nu) D_k q^0 D_i \varphi \ dx, \\ 
     I_2(j, \lambda) &\to I_2(\lambda)  = \int_{B_1^- \setminus {S_{\lambda}}} \sum_{i, k = 1}^n  F^2_{p_ip_k}(\nu) D_k q^0 D_i \varphi \ dx,
\end{align*}
as $j\to \infty$. Since $\|\nabla q^j\|_{L^2(K)}\leq C$ for $K \subset \subset B_1$, it follows that 
$$
|I_3(j, \lambda)| \leq C |S_{\lambda}|, 
$$
where the constant does not depend on $j$. If we let $\lambda\to 0$ it follows that $|I_3| \to 0$ and
$$
I_1(\lambda)+I_2(\lambda)+I_3(\lambda) \to \int_{B_1} \sum_{i, k = 1}^n  F^*_{p_ip_k}(\nu,x) D_k q^0 D_i \varphi \ dx.
$$
We therefore conclude that
\begin{equation}\label{eq:num1}
  \int_{B_1}  \sum_{i, k = 1}^n  A^j_{ik}(x) D_k q^j D_i \varphi \ dx  \to \int_{B_1} \sum_{i, k = 1}^n  F^*_{p_ip_k}(\nu,x) D_k q^0 D_i \varphi \ dx,   
\end{equation}
as $j \to \infty$. 
We next turn to the second and third integrals in \eqref{eq:pdeforqj}, and use Lemma \ref{lem:conv_AjikCmjikl} and \eqref{eq:qjass1} to conclude that
$$
\frac{|E^{m,j}_{ikl}|}{\epsilon_j} \leq C \delta_j  \|q^j\|_{C^{1, \alpha}(B_{1/2})} \leq C\|\epsilon_j q^j\|_{C^{1, \alpha}(B_{1/2})} \to 0,
$$
as $j \to \infty$ for $i,k,l=1,\hdots n$ and $m=1,2$. Thus, we see that
\begin{equation}\label{eq:num2}
 \frac{1}{\epsilon_j}  \int_{B_1^+ \setminus D^j_+}  \sum_{i, k,l = 1}^n E^{1,j}_{ikl} \nu^j_l  D^2_{ik}p^j \varphi dx +\frac{1}{\epsilon_j}   \int_{B_1^- \setminus D^j_-}  \sum_{i, k,l = 1}^n E^{2,j}_{ikl} \nu^j_l  D^2_{ik}p^j \varphi dx  \to 0,  
\end{equation}
as $j \to \infty$. For the fourth and fifth integrals in \eqref{eq:pdeforqj}, we note that $C^{1,j}_{ikl}(x)$ converges locally uniformly by Lemma \ref{lem:conv_AjikCmjikl}, and $p^j \to p^0$ in $C^{\alpha}_{\textup{loc}}(B_1)$, and therefore, since $|D^j_{\pm}| \to 0$, it follows that
 $$
 \frac{\delta_j}{\epsilon_j} \sum_{i, k, l = 1}^n \left(\int_{(B^+_1 \setminus D^j_+) \cup D^j_-}   C^{1,j}_{ikl}(x) D^2_{ik} p^j D_l p^j \varphi \ dx +\int_{(B^-_1 \setminus D^j_-) \cup D^j_+ }  C^{2,j}_{ikl}(x) D^2_{ik} p^j D_l p^j \varphi \ dx \right)
 $$
\begin{equation}\label{eq:num3}
   \to \tau \int_{B_1} \sum_{i, k, l = 1}^n   F^*_{p_ip_k p_l}(\nu,x) D^2_{ik} p^0 D_l p^0 \varphi \ dx,  
\end{equation}
as $j \to \infty$. 
We need to calculate
$$
\frac{1}{\epsilon_j} \int_{D^j_-}  \sum_{i, k = 1}^n   F^1_{p_i p_k}(\nu^j) D^2_{ik} p^j \varphi \ dx + \frac{1}{\epsilon_j} \int_{D^j_+}  \sum_{i, k = 1}^n   F^2_{p_i p_k}(\nu^j) D^2_{ik} p^j \varphi \ dx,
$$
as $j \to  \infty$. We recall that we may write $D^j_+$ and $D^j_-$ as
\begin{align*}
    D^j_- &= \{x\in B_1 :0 \geq x_n \geq - \epsilon_j f^j_-(x') + \epsilon_j g^j_-(x') \}, \\
    D^j_+ &= \{x \in B_1 :0 \leq x_n \leq -  \epsilon_j f^j_+(x') + \epsilon_j g^j_+(x') \}, 
\end{align*}
where $|g_{\pm}^j(x')| \to 0$ uniformly in $B_1$ as $j \to \infty$, and 
$$
f^j_{\pm}(x') = \frac{\frac{|\nu^j|^2-1}{\delta_j \epsilon_j} + \tau_j|\nabla p^j(x',0)|^2}{2 c^{\pm}_{\nu^j}}  \to \frac{\alpha + \tau|\nabla p^0(x',0)|^2}{2D_{\nu n} p^0_{\pm}},
$$
uniformly in $B_1'$, as $j \to \infty$. Here we have used point 2 in Lemma \ref{lem:structure_p} to see that $c_{\nu^j}^{ \pm} \to D^2_{\nu n}p^0_{\pm} \geq c_{ndg}$. This implies that, for some constant $C$, $D^j_- \subset \{x;\; |x_n|<C\epsilon_j\}$. However, it is more convenient to express $D^j_-$ and $D^j_+$ as 
\begin{align*}
    D^j_- &= \{x\in B_1 :0 \geq x_n \geq - h^j_-(x') \}, \\
    D^j_+ &= \{x \in B_1 :0 \leq x_n \leq - h^j_+(x') \},
\end{align*}
where
\begin{equation*}\label{eq:hjtoh0}
    h^j_{\pm} =  \epsilon_j f^j_-(x') - \epsilon_j g^j_-(x'),
\end{equation*}
and 
$$
\frac{h^j_{\pm}(x')}{\epsilon_j} \to \frac{\alpha + \tau |\nabla p^0(x',0)|^2}{2 D^2_{\nu n} p^0_{\pm}} \eqqcolon h^0_{\pm}(x'), 
$$ 
uniformly in $B_1'$ as $j \to \infty$. Furthermore,
$$
        \sum_{i, k = 1}^n F^1_{p_i p_k}(\nu^j)D^2_{ik} p_-^j  \to \sum_{i, k = 1}^n F^1_{p_i p_k}(\nu)D^2_{ik} p_-^0 = (\nu \cdot (D^2 F^1 (\nu)-D^2 F^2(\nu)) \nu^T) \nu_n D^2_{\nu n}p^0_-,
 $$  
 and
 $$
         \sum_{i, k = 1}^n F^2_{p_i p_k}(\nu^j)D^2_{ik} p_+^j  \to  \sum_{i, k = 1}^n F^2_{p_i p_k}(\nu)D^2_{ik} p_+^0  = -(\nu \cdot (D^2 F^1 (\nu)-D^2 F^2(\nu)) \nu^T) \nu_n D^2_{\nu n}p^0_+,
$$
by point 5 in Lemma \ref{lem:structure_p}. We recall from \eqref{eq:qjass3} that $|\nu_n| > 0$ and
\begin{equation*}
     \left| \nu \cdot (D^2 F^1(\nu)- D^2 F^2( \nu)) \nu^T \right| > 0.
\end{equation*}
In order to prove that $q^0$ satisfies the PDE \eqref{eq:C3est} we need to prove that
\begin{equation}\label{eq:num4}
    \lim_{j\to \infty} \frac{1}{\epsilon_j} \int_{D^j_-}  \varphi \ dx  =  \int_{B_1'}  \varphi(x',0) \left(\frac{\alpha+ \tau|\nabla p^0(x',0)|^2}{2 D^2_{\nu n} p^0_-}\right)^+  dx', 
\end{equation}
and
\begin{equation}\label{eq:num5}
  \lim_{j\to \infty} \frac{1}{\epsilon_j}  \int_{D^j_+}  \varphi \ dx =-  \int_{B_1'}  \varphi(x',0) \left(\frac{\alpha+ \tau|\nabla p^0(x',0)|^2}{2 D^2_{\nu n} p^0_+}\right)^-  dx'.  
\end{equation}
We will focus on the integral over $D^j_-$ and note that the integral over $D^j_+$ can be calculated using the same method. Next, we define the functional $F^j_-: C^1_0(B_1)\to \mathbb{R}$ by
$$
F^j_-(\varphi)=\frac{1}{\epsilon_j}\int_{D^j_-}  \varphi(x) \ dx_n dx' = \frac{1}{\epsilon_j}\int_{B_1'} \int_{-(h^j_-)^+}^0  \varphi(x) \ dx'.
$$

We choose $\gamma>0$ and $j$ so large that $|\varphi(x',0)-\varphi(x)| \leq \gamma$, when $x_n \in (-(h^j_-)^+),0)$. Then
\begin{align*}
\left|F^j_-(\varphi)-\frac{1}{\epsilon_j} \int_{B_1'} \int_{-(h^j_-)^+}^0 \varphi(x',0) dx_n dx' \right| &\leq \frac{\gamma}{\epsilon_j} \int_{B_1'} \int_{-(h^j_-)^+}^0 dx_n dx' \\
&= \int_{B_1'} \frac{\gamma(h^j_-)^+}{\epsilon_j} dx'\\
&\leq C \gamma,
\end{align*}
where we used \eqref{eq:hjtoh0} in the final inequality. Since
$$
\frac{1}{\epsilon_j} \int_{B_1'} \int_{-(h^j_-)^+}^0 \varphi(x',0) dx_n dx' = \int_{B_1'} \frac{(h^j_-)^+}{\epsilon_j}  \varphi(x',0) dx' \to \int_{B_1'}  (h^0_-(x'))^+ \varphi(x',0) dx',
$$
as $j \to \infty$ we have proved \eqref{eq:num4}.

Next, we let $j$ tend to infinity, and use \eqref{eq:num1}-\eqref{eq:num5} which proves \eqref{eq:q0equation}. The estimates in \eqref{eq:C3est} now follows directly from Corollary \ref{cor:C3result} with
$$
m(x)=  \frac{\nu \cdot (D^2 F^1 (\nu)-D^2 F^2(\nu)) \nu^T }{2}\nu_n\tau|\nabla p^0(x',0)|^2,
$$
$$
M =  \frac{\nu \cdot (D^2 F^1 (\nu)-D^2 F^2(\nu))\nu^T}{2} \nu_n \alpha,
$$
and
 $$
 r(x)=\left\{ \begin{array}{ll}
               \tau \sum_{i, k, l = 1}^n   F^1_{p_ip_k p_l}(\nu) D_{ik} p^0 D_l p^0  & \textrm{for }x_n\ge 0, \\
               \tau \sum_{i, k, l = 1}^n   F^2_{p_ip_k p_l}(\nu) D_{ik} p^0 D_l p^0 & \textrm{for }x_n<0,
              \end{array}
 \right.
 $$
which clearly is linear in $B_1^+$ and $B_1^-$, since $p^0$ is a homogeneous second order polynomial in $B_1^+$ and $B_1^-$, respectively. Thus, the theorem is proved.
\end{proof}

\end{document}
